\VerdierPhysicalPageAnchor{0042}%
\VerdierSemanticAnchor{ch1-1-3-7-cont}
exclusively in this work the language of $G^\circ$-categories. However,
we shall not entirely forget the \textquotedblleft graded
categories\textquotedblright\ viewpoint, using the following
terminology:

\VerdierSemanticAnchor{ch1-1-3-8}
\noindent\textbf{Definition 1.3.8.}\quad
Let $G$ be a group and let $(\mathcal C,G^\circ,T,c)$ be a
$G^\circ$-category. Let $g$ and $h$ be two elements of $G$, and let
$X$ and $Y$ be two objects of $\mathcal C$. A morphism in $\mathcal C$

\VerdierSemanticAnchor{ch1-1-3-8-morphism}
\[
  u\colon T(h)X\longrightarrow T(gh)Y
\]
will be called \emph{a morphism of degree $g$ from $X$ to $Y$}.

Let us show that this terminology causes no confusion. Indeed, for
every pair $h,h'$ of elements of $G$, the functor $T(h^{-1}h')$ and
the isomorphisms $c(h,h^{-1}h')$ and $c(gh,h^{-1}h')$ define a
bifunctorial isomorphism

\VerdierSemanticAnchor{ch1-1-3-8-eq1}
\begin{equation}
  i_g(h,h')\colon
  \operatorname{Hom}_{\mathcal C}\bigl(T(h)X,T(gh)Y\bigr)
  \overset{\sim}{\longrightarrow}
  \operatorname{Hom}_{\mathcal C}\bigl(T(h')X,T(gh')Y\bigr).
  \tag{1.3.8.1}
\end{equation}

The cocycle identities (1.2.1.3) show that one has the relation

\VerdierSemanticAnchor{ch1-1-3-8-eq2}
\begin{equation}
  i_g(h,h'')=i_g(h',h'')i_g(h,h'),
  \tag{1.3.8.2}
\end{equation}
for every triple $h,h',h''\in G$. Thus, if the different sets

\VerdierSemanticAnchor{ch1-1-3-8-identification}
\[
  \operatorname{Hom}_{\mathcal C}\bigl(T(h)X,T(gh)Y\bigr)
\]
are identified by means of the isomorphisms $i_g(h,h')$, formula
(1.3.8.2) shows that these different identifications are compatible.
Moreover, for every quadruple $g,g',h,h'\in G$ and every triple
$X,Y,Z$ of objects of $\mathcal C$, diagram (1.3.8.3) below commutes:

\VerdierSemanticAnchor{ch1-1-3-8-eq3}
\begin{center}
  \resizebox{0.92\textwidth}{!}{%
    $\begin{tikzcd}[ampersand replacement=\&,column sep=5em,row sep=4em]
      \operatorname{Hom}_{\mathcal C}\bigl(T(h)X,T(gh)Y\bigr)
      \times
      \operatorname{Hom}_{\mathcal C}\bigl(T(gh)Y,T(g'gh)Z\bigr)
      \arrow[r]
      \arrow[d,
        "{i_g(h,h')\times i_{g'}(gh,gh')}"',
        "\sim"]
      \&
      \operatorname{Hom}_{\mathcal C}\bigl(T(h)X,T(g'gh)Z\bigr)
      \arrow[d,"{i_{g'g}(h,h')}"',"\sim"]
      \\
      \operatorname{Hom}_{\mathcal C}\bigl(T(h')X,T(gh')Y\bigr)
      \times
      \operatorname{Hom}_{\mathcal C}\bigl(T(gh')Y,T(g'gh')Z\bigr)
      \arrow[r]
      \&
      \operatorname{Hom}_{\mathcal C}\bigl(T(h')X,T(g'gh')Z\bigr)
    \end{tikzcd}$%
  }
  \hfill\textup{(1.3.8.3)}
\end{center}

\VerdierSemanticAnchor{ch1-1-3-8-composition}
(the horizontal arrows are defined by composition in $\mathcal C$).
Diagram (1.3.8.3) therefore allows us to define without ambiguity, or
more precisely up to the identifications defined by the isomorphisms
$i_g(h,h')$, the composite of two morphisms of respective degrees $g$
and $g'$. One verifies immediately that this definition of composition
is compatible with the definition
